\documentclass[journal, 12pt, onecolumn]{IEEEtran}
\usepackage[latin1]{inputenc}
\usepackage[spanish]{babel}
\usepackage{amsfonts}
\usepackage{amsmath}
\usepackage[dvips]{graphicx}

\begin{document}
\title{\textit{Informe Entrega 2 - Paralelo Y - Grupo Z}}
\author{Nombres integrantes \\ 
Correos integrantes \\
UTFSM Casa Central, Valpara\'iso, Chile}
\maketitle

\section{Introducci\'on}
Describa en qu\'e consiste el problema, la necesidad que lo motiva, el contexto de
aplicaci\'on y el objetivo general del proyecto. Incluya los supuestos principales
adoptados para representar la situaci\'on real. (5 puntos)

\section{Modelamiento del problema}
Formule expl\'icitamente el modelo matem\'atico utilizando la materia vista en
clases. El documento debe incluir, como m\'inimo:
\begin{itemize}
    \item los conjuntos e \textit{\'indices} utilizados;
    \item los par\'ametros, sus unidades, fuentes y rangos de valores;
    \item las variables de decisi\'on y sus dominios (por ejemplo, binarias,
    enteras o continuas);
    \item la funci\'on objetivo, indicando si se minimiza o maximiza, y
    explicando el significado de cada t\'ermino;
    \item todas las restricciones del modelo, acompa\~nadas de una breve
    explicaci\'on de su significado f\'isico u operativo;
    \item las condiciones de dominio y cualquier supuesto adicional necesario.
\end{itemize}
(5 puntos)

\section{Generador de Instancias}
Describa detalladamente el procedimiento de generaci\'on de las \textbf{15
instancias factibles} solicitadas. Explique:
\begin{itemize}
	\item el algoritmo o rutina implementada, incluyendo sus entradas, salidas y
	par\'ametros configurables;
	\item las distribuciones, semillas, rangos y dem\'as decisiones usadas para
	producir aleatoriedad reproducible;
	\item el tama\~no y la estructura de cada instancia, y la forma en que se
	almacenan los datos;
    \item el procedimiento de validaci\'on que comprueba que cada instancia
    satisface todas las restricciones;
    \item las condiciones implementadas para garantizar la factibilidad
    t\'ecnica de las soluciones generadas.
\end{itemize}
Incluya una tabla resumen de las 15 instancias y entregue el c\'odigo completo
del generador. (5 puntos)

\section{Implementaci\'on}
Explique el solver asignado y la implementaci\'on realizada (por ejemplo,
lp\_solve o MiniZinc), incluyendo:
\begin{itemize}
	\item como el solver realiza la b\'usqueda de la soluci\'on \textbf{\'optima};
	\item la construcci\'on automatizada del modelo a partir de los datos de
	entrada;
	\item el procedimiento para ejecutar y resolver cada instancia;
	\item los par\'ametros del solver y el criterio de optimalidad utilizado;
	\item el formato de las soluciones retornadas y las comprobaciones de
	factibilidad y optimalidad.
\end{itemize}
El c\'odigo debe poder parametrizar los datos de entrada, procesar las
instancias generadas y retornar la soluci\'on \textbf{\'optima} de cada una.
Entregue el c\'odigo en un archivo independiente.
(30 puntos)

\section{Resultados}
Presente un an\'alisis estructurado de los experimentos computacionales sobre
las 15 instancias, respaldado por gr\'aficos y tablas. Reporte, como m\'inimo,
los valores de la soluci\'on \textbf{\'optima}, el estado de factibilidad, el
tiempo de ejecuci\'on y las m\'etricas relevantes del solver.

\subsection{Impacto de los m\'argenes relativos}
Analice c\'omo afectan los m\'argenes relativos del m\'aximo permitido para la
demanda energ\'etica y del costo total a la factibilidad, al valor objetivo y al
espacio de b\'usqueda. Compare escenarios y explique los cambios observados.

\subsection{Desempe\~no computacional}
Eval\'ue la dimensionalidad de cada instancia, indicando el n\'umero de
variables binarias y de restricciones, frente al tiempo de ejecuci\'on, los
nodos explorados y la brecha de optimalidad (\textit{MIP gap}), cuando estas
m\'etricas est\'en disponibles. Discuta tendencias, casos extremos y posibles
causas de las diferencias de desempe\~no. (45 puntos)

\section{Conclusiones}
Actualice las conclusiones de la Entrega 1 a partir de los resultados de esta
entrega. Resuma los hallazgos sobre el modelo, las instancias y la
implementaci\'on del solver, indicando la calidad y optimalidad de las
soluciones obtenidas. Analice el efecto de los m\'argenes relativos sobre la
factibilidad y el espacio de b\'usqueda, y relacione la dimensionalidad de las
instancias con el tiempo de ejecuci\'on, los nodos explorados y la brecha de
optimalidad. (10 puntos)

\end{document}
